% Options for packages loaded elsewhere
\PassOptionsToPackage{unicode}{hyperref}
\PassOptionsToPackage{hyphens}{url}
%
\documentclass[
]{article}
\usepackage{amsmath,amssymb}
\usepackage{iftex}
\ifPDFTeX
  \usepackage[T1]{fontenc}
  \usepackage[utf8]{inputenc}
  \usepackage{textcomp} % provide euro and other symbols
\else % if luatex or xetex
  \usepackage{unicode-math} % this also loads fontspec
  \defaultfontfeatures{Scale=MatchLowercase}
  \defaultfontfeatures[\rmfamily]{Ligatures=TeX,Scale=1}
\fi
\usepackage{lmodern}
\ifPDFTeX\else
  % xetex/luatex font selection
\fi
% Use upquote if available, for straight quotes in verbatim environments
\IfFileExists{upquote.sty}{\usepackage{upquote}}{}
\IfFileExists{microtype.sty}{% use microtype if available
  \usepackage[]{microtype}
  \UseMicrotypeSet[protrusion]{basicmath} % disable protrusion for tt fonts
}{}
\makeatletter
\@ifundefined{KOMAClassName}{% if non-KOMA class
  \IfFileExists{parskip.sty}{%
    \usepackage{parskip}
  }{% else
    \setlength{\parindent}{0pt}
    \setlength{\parskip}{6pt plus 2pt minus 1pt}}
}{% if KOMA class
  \KOMAoptions{parskip=half}}
\makeatother
\usepackage{xcolor}
\usepackage{color}
\usepackage{fancyvrb}
\newcommand{\VerbBar}{|}
\newcommand{\VERB}{\Verb[commandchars=\\\{\}]}
\DefineVerbatimEnvironment{Highlighting}{Verbatim}{commandchars=\\\{\}}
% Add ',fontsize=\small' for more characters per line
\newenvironment{Shaded}{}{}
\newcommand{\AlertTok}[1]{\textcolor[rgb]{1.00,0.00,0.00}{\textbf{#1}}}
\newcommand{\AnnotationTok}[1]{\textcolor[rgb]{0.38,0.63,0.69}{\textbf{\textit{#1}}}}
\newcommand{\AttributeTok}[1]{\textcolor[rgb]{0.49,0.56,0.16}{#1}}
\newcommand{\BaseNTok}[1]{\textcolor[rgb]{0.25,0.63,0.44}{#1}}
\newcommand{\BuiltInTok}[1]{\textcolor[rgb]{0.00,0.50,0.00}{#1}}
\newcommand{\CharTok}[1]{\textcolor[rgb]{0.25,0.44,0.63}{#1}}
\newcommand{\CommentTok}[1]{\textcolor[rgb]{0.38,0.63,0.69}{\textit{#1}}}
\newcommand{\CommentVarTok}[1]{\textcolor[rgb]{0.38,0.63,0.69}{\textbf{\textit{#1}}}}
\newcommand{\ConstantTok}[1]{\textcolor[rgb]{0.53,0.00,0.00}{#1}}
\newcommand{\ControlFlowTok}[1]{\textcolor[rgb]{0.00,0.44,0.13}{\textbf{#1}}}
\newcommand{\DataTypeTok}[1]{\textcolor[rgb]{0.56,0.13,0.00}{#1}}
\newcommand{\DecValTok}[1]{\textcolor[rgb]{0.25,0.63,0.44}{#1}}
\newcommand{\DocumentationTok}[1]{\textcolor[rgb]{0.73,0.13,0.13}{\textit{#1}}}
\newcommand{\ErrorTok}[1]{\textcolor[rgb]{1.00,0.00,0.00}{\textbf{#1}}}
\newcommand{\ExtensionTok}[1]{#1}
\newcommand{\FloatTok}[1]{\textcolor[rgb]{0.25,0.63,0.44}{#1}}
\newcommand{\FunctionTok}[1]{\textcolor[rgb]{0.02,0.16,0.49}{#1}}
\newcommand{\ImportTok}[1]{\textcolor[rgb]{0.00,0.50,0.00}{\textbf{#1}}}
\newcommand{\InformationTok}[1]{\textcolor[rgb]{0.38,0.63,0.69}{\textbf{\textit{#1}}}}
\newcommand{\KeywordTok}[1]{\textcolor[rgb]{0.00,0.44,0.13}{\textbf{#1}}}
\newcommand{\NormalTok}[1]{#1}
\newcommand{\OperatorTok}[1]{\textcolor[rgb]{0.40,0.40,0.40}{#1}}
\newcommand{\OtherTok}[1]{\textcolor[rgb]{0.00,0.44,0.13}{#1}}
\newcommand{\PreprocessorTok}[1]{\textcolor[rgb]{0.74,0.48,0.00}{#1}}
\newcommand{\RegionMarkerTok}[1]{#1}
\newcommand{\SpecialCharTok}[1]{\textcolor[rgb]{0.25,0.44,0.63}{#1}}
\newcommand{\SpecialStringTok}[1]{\textcolor[rgb]{0.73,0.40,0.53}{#1}}
\newcommand{\StringTok}[1]{\textcolor[rgb]{0.25,0.44,0.63}{#1}}
\newcommand{\VariableTok}[1]{\textcolor[rgb]{0.10,0.09,0.49}{#1}}
\newcommand{\VerbatimStringTok}[1]{\textcolor[rgb]{0.25,0.44,0.63}{#1}}
\newcommand{\WarningTok}[1]{\textcolor[rgb]{0.38,0.63,0.69}{\textbf{\textit{#1}}}}
\usepackage{longtable,booktabs,array}
\usepackage{calc} % for calculating minipage widths
% Correct order of tables after \paragraph or \subparagraph
\usepackage{etoolbox}
\makeatletter
\patchcmd\longtable{\par}{\if@noskipsec\mbox{}\fi\par}{}{}
\makeatother
% Allow footnotes in longtable head/foot
\IfFileExists{footnotehyper.sty}{\usepackage{footnotehyper}}{\usepackage{footnote}}
\makesavenoteenv{longtable}
\setlength{\emergencystretch}{3em} % prevent overfull lines
\providecommand{\tightlist}{%
  \setlength{\itemsep}{0pt}\setlength{\parskip}{0pt}}
\setcounter{secnumdepth}{-\maxdimen} % remove section numbering
% definitions for citeproc citations
\NewDocumentCommand\citeproctext{}{}
\NewDocumentCommand\citeproc{mm}{%
  \begingroup\def\citeproctext{#2}\cite{#1}\endgroup}
\makeatletter
 % allow citations to break across lines
 \let\@cite@ofmt\@firstofone
 % avoid brackets around text for \cite:
 \def\@biblabel#1{}
 \def\@cite#1#2{{#1\if@tempswa , #2\fi}}
\makeatother
\newlength{\cslhangindent}
\setlength{\cslhangindent}{1.5em}
\newlength{\csllabelwidth}
\setlength{\csllabelwidth}{3em}
\newenvironment{CSLReferences}[2] % #1 hanging-indent, #2 entry-spacing
 {\begin{list}{}{%
  \setlength{\itemindent}{0pt}
  \setlength{\leftmargin}{0pt}
  \setlength{\parsep}{0pt}
  % turn on hanging indent if param 1 is 1
  \ifodd #1
   \setlength{\leftmargin}{\cslhangindent}
   \setlength{\itemindent}{-1\cslhangindent}
  \fi
  % set entry spacing
  \setlength{\itemsep}{#2\baselineskip}}}
 {\end{list}}
\usepackage{calc}
\newcommand{\CSLBlock}[1]{\hfill\break\parbox[t]{\linewidth}{\strut\ignorespaces#1\strut}}
\newcommand{\CSLLeftMargin}[1]{\parbox[t]{\csllabelwidth}{\strut#1\strut}}
\newcommand{\CSLRightInline}[1]{\parbox[t]{\linewidth - \csllabelwidth}{\strut#1\strut}}
\newcommand{\CSLIndent}[1]{\hspace{\cslhangindent}#1}
\ifLuaTeX
  \usepackage{selnolig}  % disable illegal ligatures
\fi
\usepackage{bookmark}
\IfFileExists{xurl.sty}{\usepackage{xurl}}{} % add URL line breaks if available
\urlstyle{same}
\hypersetup{
  pdftitle={{[}Re{]} Computational Replication of SADDLE: Empirical Next-Generation Sequencing Validation of Multiplex PCR Read Partitions},
  hidelinks,
  pdfcreator={LaTeX via pandoc}}

\title{{[}Re{]} Computational Replication of SADDLE: Empirical
Next-Generation Sequencing Validation of Multiplex PCR Read Partitions}
\author{true}
\date{2026-11-10}

\begin{document}
\maketitle

\section{Abstract}\label{abstract}

Multiplex polymerase chain reaction (multiplex PCR) is a foundational
enrichment technology for targeted next-generation sequencing (NGS),
clinical pathogen panels, and forensic identification. However,
high-plex reactions suffer from non-specific primer-dimer formation,
which can monopolize sequencing throughput and deplete on-target
amplification. In 2022, Xie et al.~introduced SADDLE (\emph{Simulated
Annealing Design using Dimer Likelihood Estimation}, \emph{Nature
Communications}, 13, 1881; PMID: 35410464), describing an empirical
heuristic scoring framework and simulated annealing algorithm to design
multiplex primer pools up to 384-plex. While the experimental utility of
SADDLE was demonstrated in the original paper, the validation sequencing
analysis relied on proprietary MATLAB scripts coupled to external
command-line aligners, while the simulated annealing design code
remained proprietary under patent restrictions. In this work, following
the ReScience C guidelines for computational replication, we conduct an
independent replication of the central empirical NGS validation finding
of Xie et al.: the read classification breakdown and primer-interaction
distribution of the SADDLE-designed 384-plex validation panel across
3,942,385 raw paired-end sequencing reads (Sample E19; Figure 3f). Using
an independent, open-source Python reimplementation, our pipeline
classifies 2,860,658 reads into 43.86\% on-target amplicons, 55.27\%
nonspecific products, and 0.87\% primer-dimers, reproducing the
published proportions (\textasciitilde43\%, \textasciitilde56\%,
\textasciitilde1\%) within a 1.0\% absolute tolerance. We further
corroborate the severe heavy-tailed concentration of primer-dimer
formation across the 295,296 candidate primer pairings, where only 1,172
pairs (0.397\%) exhibit non-zero dimerization and a single interaction
(\texttt{rp213} \(\times\) \texttt{rp371}) accounts for 67.57\% of all
dimer reads (16,805 reads). In a supplementary exploratory benchmark, we
evaluate the predictive validity of the SADDLE Badness heuristic against
nearest-neighbor thermodynamics and supervised machine learning models,
documenting an Average Precision of 0.0597 (a 4.8-fold improvement over
nearest-neighbor free energy alone) while analyzing the severe
probability miscalibration of supervised baselines (Brier scores
\textasciitilde0.20 vs \textasciitilde0.0040 baseline). All replication
code, derived datasets, and an automated verification script that
executes in under two seconds are openly accessible under permissive
licenses.

\section{Introduction}\label{introduction}

Multiplex PCR simultaneously amplifies dozens to hundreds of genomic
loci within a single reaction volume, serving as an essential
preparation step for targeted clinical diagnostics, comprehensive cancer
hotspot panels, and metagenomic surveillance (Untergasser et al. 2012;
Chen et al. 2026). However, scaling multiplex PCR beyond small panels is
fundamentally constrained by pairwise primer-primer cross-hybridization.
For a panel targeting \(N\) amplicons, \(2N\) distinct primer
oligonucleotides are pooled, yielding:

\[\binom{2N+1}{2} = \frac{(2N+1)(2N)}{2} \approx 2N^2\]

possible pairwise interaction combinations (including homodimer and
self-interactions). Minor sequence complementarity---specifically when
aligned at the extensible 3' terminus---enables DNA polymerase
extension, producing short artifactual amplicons termed primer-dimers
(Xie et al. 2022). Due to their substantially smaller size relative to
target genomic amplicons, primer-dimers amplify with high kinetic
efficiency, rapidly exhausting enzymatic reagents and saturating
sequencing libraries.

To mitigate this exponential risk, Xie et al. (2022) introduced SADDLE
(\textbf{S}imulated \textbf{A}nnealing \textbf{D}esign using
\textbf{D}imer \textbf{L}ikelihood \textbf{E}stimation), proposing an
empirical biophysical scoring heuristic called \emph{Badness} to
quantify 3'-directed extensible hybridization and using simulated
annealing optimization to partition multiplex primer pools into
compatible sub-pools. In their 96-plex experimental optimization
(Figures 3a--3c), the authors demonstrated that primer-dimer reads
dropped from 90.7\% in a naively designed pool (PS1) to 4.9\% in an
optimized pool (PS3). Scaling the technology to higher multiplexing, the
authors designed and experimentally tested a \textbf{384-plex primer set
designed by SADDLE (768 primers; Sample E19)}. In \textbf{Figure 3f} and
the accompanying text, the authors reported that the sequencing library
of this 384-plex panel achieved high on-target specificity
(\textasciitilde43\%) and low primer-dimer formation
(\textasciitilde1\%), with the remaining reads (\textasciitilde56\%)
representing nonspecific background products aligning across the human
genome via non-cognate primer pairs.

Despite the impact of SADDLE, two key reproducibility hurdles exist: 1.
\textbf{Proprietary Software Barriers:} The simulated annealing software
used to design primer sets was withheld under intellectual property and
patent restrictions (Xie et al. 2022). Furthermore, the empirical NGS
validation scripts deposited on Figshare were implemented in proprietary
MATLAB and required external aligners (Langmead and Salzberg 2012),
preventing turnkey, open-source execution. 2. \textbf{Independent
Verification of Empirical Partitions:} While the authors published
rounded percentage summaries in Figure 3f of their manuscript, an
independent, open-source replication of the read-accounting pipeline
operating directly on the raw FASTQ data has not been formally
documented.

In this ReScience C article (Rougier et al. 2017), we report an
independent computational replication focused strictly on the empirical
validation analysis of Xie et al. (2022): - \textbf{Primary Replication
Objective:} Does an independent, open-source Python implementation
replicate the published read classification proportions
(\textasciitilde43\% on-target, \textasciitilde56\% nonspecific,
\textasciitilde1\% primer-dimer) across all 3,942,385 raw paired-end
reads of the SADDLE 384-plex validation panel (Sample E19, Figure 3f),
and does it reproduce the heavy-tailed concentration of primer-dimer
reads? - \textbf{Implementation Mapping:} Can the read-classification
pipeline be reconstructed without proprietary MATLAB dependencies or
external binary alignment tools while maintaining operational fidelity?
- \textbf{Supplementary Exploratory Analysis:} In a secondary analysis
clearly distinguished from the primary replication claim, how well does
the SADDLE Badness heuristic rank observed dimer events compared to
standard nearest-neighbor thermodynamics (\(-\Delta G\)) and supervised
machine learning models?

\section{Original Work and Artifact
Audit}\label{original-work-and-artifact-audit}

The original study by Xie et al. (2022) was published in \emph{Nature
Communications} (Volume 13, Article 1881; PMID: 35410464). Primer
sequences, genomic coordinates, and assay details were provided in
Supplementary Data (MOESM5). The raw paired-end sequencing reads for the
384-plex validation panel (Sample E19, Illumina MiniSeq 2 \(\times\) 150
bp) and post-hoc analysis scripts (\texttt{Data\_process.m},
\texttt{Dimer\_Identify.m}, \texttt{analysis\_code.ipynb}) were
deposited on Figshare (DOI: 10.6084/m9.figshare.16944154.v3) under a
Creative Commons Attribution 4.0 International license (CC BY 4.0).

Our audit of the deposited artifacts identified: - \textbf{Design Code:}
The simulated annealing optimization algorithm was not included in the
public deposit, with the manuscript stating: \emph{``The code for SADDLE
primer design is available upon request for non-commercial academic
use''} subject to patent protection. Consequently, the design algorithm
itself is out of scope for an open, autonomous replication. -
\textbf{Validation Pipeline:} The public analysis workflow comprises
MATLAB scripts that trim Illumina adapters, filter reads by Phred
quality and length, align reads to reference amplicons using Bowtie2
(Langmead and Salzberg 2012), match 5' primer sequences via an 8-base
hash lookup, and classify reads with insert lengths strictly shorter
than \(L_1 + L_2 - 2\) bp as primer-dimers. - \textbf{Raw Data
Completeness:} The Figshare repository contains the complete 3,942,385
raw paired-end reads for Sample E19 across two gzipped FASTQ files
(\texttt{E19\_384plex\_NGS\_R1\_fastq.gz} and
\texttt{E19\_384plex\_NGS\_R2\_fastq.gz}, totaling 914 MB), providing
uncompromised primary empirical inputs.

\section{Replication Methodology and
Implementation}\label{replication-methodology-and-implementation}

We implemented a clean-room, open-source Python 3.10+ pipeline designed
to operate without MATLAB or platform-dependent aligners.

\subsection{Structural Code Mapping}\label{structural-code-mapping}

Table 1 details the structural correspondence between the original
SADDLE validation code and our Python reconstruction.

\begin{longtable}[]{@{}
  >{\raggedright\arraybackslash}p{(\columnwidth - 6\tabcolsep) * \real{0.2500}}
  >{\raggedright\arraybackslash}p{(\columnwidth - 6\tabcolsep) * \real{0.2500}}
  >{\raggedright\arraybackslash}p{(\columnwidth - 6\tabcolsep) * \real{0.2500}}
  >{\raggedright\arraybackslash}p{(\columnwidth - 6\tabcolsep) * \real{0.2500}}@{}}
\caption{Structural mapping and methodological comparison between the
published SADDLE validation scripts and the clean-room Python
implementation.}\tabularnewline
\toprule\noalign{}
\begin{minipage}[b]{\linewidth}\raggedright
Category
\end{minipage} & \begin{minipage}[b]{\linewidth}\raggedright
Original SADDLE Logic (MATLAB / Bowtie2)
\end{minipage} & \begin{minipage}[b]{\linewidth}\raggedright
Python Clean-Room Reconstruction
\end{minipage} & \begin{minipage}[b]{\linewidth}\raggedright
Equivalence Classification
\end{minipage} \\
\midrule\noalign{}
\endfirsthead
\toprule\noalign{}
\begin{minipage}[b]{\linewidth}\raggedright
Category
\end{minipage} & \begin{minipage}[b]{\linewidth}\raggedright
Original SADDLE Logic (MATLAB / Bowtie2)
\end{minipage} & \begin{minipage}[b]{\linewidth}\raggedright
Python Clean-Room Reconstruction
\end{minipage} & \begin{minipage}[b]{\linewidth}\raggedright
Equivalence Classification
\end{minipage} \\
\midrule\noalign{}
\endhead
\bottomrule\noalign{}
\endlastfoot
\textbf{Adapter Trimming} & \texttt{Data\_process.m:23-57}: Scans
\texttt{AGATCGGAAGAGCACACGTCTGAACTCC} on R1 and
\texttt{AGATCGGAAGAGCGTCGTGTAGGGAAAG} on R2; trims read to prefix before
adapter. & \texttt{reconstruct\_saddle\_384.py:145-152}: Scans core
17-bp Illumina TruSeq prefixes \texttt{AGATCGGAAGAGCACAC} (R1) and
\texttt{AGATCGGAAGAGCGTCG} (R2); trims reads. & Operationally
Equivalent \\
\textbf{Quality Filtering} & \texttt{Data\_process.m:109-111}: Filters
read pairs unless \(>70\%\) of bases have Phred quality \(>30\) (ASCII
\(>63\)), no ambiguous \texttt{\textquotesingle{}N\textquotesingle{}}
bases, and trimmed length \(>10\) bp. &
\texttt{reconstruct\_saddle\_384.py:155-163}: Exact condition: length
\(>10\), \texttt{\textquotesingle{}N\textquotesingle{}} not in sequence,
and \(\ge 70\%\) bases with quality ASCII \(>63\) for both R1 and R2. &
Exact Numerical Match \\
\textbf{Primer Matching} & \texttt{Dimer\_Identify.m:28-41,\ 71-95}:
Computes 8-base hash using \texttt{PrimeNumberHash(lower(p\{i\}(1:8)))};
evaluates Smith-Waterman prefix identity \(>0.6\). &
\texttt{reconstruct\_saddle\_384.py:166-196}: Builds 8-bp prefix index
across all 768 primers; evaluates prefix match identity \(>0.6\) on
prefix of length \(\min(\text{len}(p), \text{len}(t))\). & Operationally
Equivalent \\
\textbf{Insert Length \& Dimer Threshold} &
\texttt{Dimer\_Identify.m:96-108}: Dimer condition:
\texttt{length(currline1)\ -\ 1\ \textless{}\ (length(f\_seq)\ +\ length(r\_seq)\ -\ 2)}.
& \texttt{reconstruct\_saddle\_384.py:203-220}: Dimer threshold:
\(L_1 + L_2 - 2\). Reads with insert length strictly below threshold are
classified as dimers. & Exact Formula Match \\
\textbf{On-Target Amplicon Assignment} &
\texttt{analysis\_code.ipynb:17-38}: Executes Bowtie2 alignment against
reference amplicons; unmapped reads passed to dimer script. &
\texttt{reconstruct\_saddle\_384.py:199-210}: Identifies cognate pairs
by primer index (\(P_2 = P_1 + 384\)); classifies cognate reads with
insert length \(\ge L_1 + L_2 - 2\) as on-target. & Methodological
Approximation (Primer pairing vs.~sequence alignment) \\
\textbf{Pair Aggregation \& Homodimers} &
\texttt{Dimer\_Identify.m:98,\ 125-126}: Increments \(768 \times 768\)
matrix; preserves homodimers on diagonal \texttt{dimer(i,\ i)}. &
\texttt{reconstruct\_saddle\_384.py:238-249}: Computes
\(768 \times 768\) ordered matrix, then folds into 295,296 unordered
pairs with diagonal preserved. & Exact Match \\
\end{longtable}

\subsection{Read Accounting and Partition
Logic}\label{read-accounting-and-partition-logic}

Each paired-end sequencing read \((R_1, R_2)\) in Sample E19 is
processed through mutually exclusive and exhaustive accounting bins: 1.
\textbf{Adapter Missing:} Read pairs lacking adapter signatures on both
reads, R1 only, or R2 only. 2. \textbf{Quality Excluded:} Read pairs
with trimmed inserts \(\le 10\) bp, containing ambiguous
\texttt{\textquotesingle{}N\textquotesingle{}} bases, or failing the
70\% Q30 threshold. 3. \textbf{Primer Unmatched:} Read pairs where
either or both 5' ends cannot be matched to the 768 pool primers with
\(>60\%\) identity. 4. \textbf{Classified Amplicons:} Reads with both
primers successfully resolved. These are partitioned into
\textbf{Primer-Dimers} (insert \(< L_1 + L_2 - 2\) bp), and non-dimer
products. To evaluate on-target specificity rigorously, non-dimer
products are assessed under two distinct methodological criteria: -
\textbf{Method 1 (Primer Pairing Heuristic):} A read is classified as
\emph{On-Target} if both primers form a cognate target pair
(\(P_2 = P_1 + 384\)) and insert length \(\ge L_1 + L_2 - 2\) bp; all
non-cognate pairings are classified as \emph{Nonspecific}. This turnkey
method avoids external aligner dependencies. - \textbf{Method 2 (Genuine
Gapped Sequence Alignment via edlib):} In addition to cognate primer
matching, the insert sequence is aligned directly against the intended
full-length reference amplicon sequence extracted from Supplementary
Data MOESM5 (Xie et al. 2022) using the open-source \texttt{edlib}
C/Python library (Myers' exact bit-vector semi-global alignment with
indels and substitutions,
\texttt{mode=\textquotesingle{}SHW\textquotesingle{}}). Reads achieving
\(\ge 75\%\) or \(\ge 80\%\) sequence identity against the intended
amplicon are classified as \emph{On-Target}; reads failing this identity
threshold are classified as \emph{Nonspecific}.

\section{Primary Replication Results}\label{primary-replication-results}

\subsection{Replicated NGS Read Accounting
Breakdown}\label{replicated-ngs-read-accounting-breakdown}

Executing the read-accounting pipeline across all 3,942,385 raw
paired-end reads of Sample E19 classifies 2,860,658 to 2,860,664 reads
(72.56\% of total raw sequencing depth), with 1,081,721 to 1,081,727
reads (27.44\%) excluded. Exclusions partition into: - Missing adapter
sequence: 751,644 reads (19.07\% of raw depth); - Quality and length
filters: 137,086 reads (3.48\% of raw depth); - Primer prefix unmatched:
192,991 to 192,997 reads (4.90\% of raw depth).

\emph{(Note on the 6-Read Lineage Discrepancy: The minor difference of
exactly 6 reads between accounting partitions arises from the
denominator used during primer prefix matching. In extreme-short primer
dimers with inserts of 12--14 bp, dividing matches by full primer length
(\(\approx 22\) bp) yields identity \(\le 0.60\) and excludes the reads
as primer-unmatched {[}192,997 excluded, 24,872 dimers{]}, whereas
dividing by the trimmed insert length evaluates identity as \(1.0\) and
retains them as dimers {[}192,991 excluded, 24,878 dimers{]}. This
single-read lineage difference does not affect any published
percentage).}

Table 2 compares the classified read proportions directly against the
published findings in Figure 3f and the main text of Xie et al. (2022)
across all evaluation methodologies.

\begin{longtable}[]{@{}
  >{\raggedright\arraybackslash}p{(\columnwidth - 10\tabcolsep) * \real{0.1667}}
  >{\raggedright\arraybackslash}p{(\columnwidth - 10\tabcolsep) * \real{0.1667}}
  >{\raggedright\arraybackslash}p{(\columnwidth - 10\tabcolsep) * \real{0.1667}}
  >{\raggedright\arraybackslash}p{(\columnwidth - 10\tabcolsep) * \real{0.1667}}
  >{\raggedright\arraybackslash}p{(\columnwidth - 10\tabcolsep) * \real{0.1667}}
  >{\raggedright\arraybackslash}p{(\columnwidth - 10\tabcolsep) * \real{0.1667}}@{}}
\caption{Methodological comparison of replicated read classification
proportions against published SADDLE 384-plex results (Sample E19,
MiniSeq 2 \(\times\) 150 bp, Figure 3f). Method 1 evaluates cognate
primer pairing; Method 2 evaluates genuine gapped semi-global sequence
alignment via \texttt{edlib} against all 384 intended reference
amplicons.}\tabularnewline
\toprule\noalign{}
\begin{minipage}[b]{\linewidth}\raggedright
Read Category
\end{minipage} & \begin{minipage}[b]{\linewidth}\raggedright
Method 1: Primer Pairing Count (\%)
\end{minipage} & \begin{minipage}[b]{\linewidth}\raggedright
Method 2: edlib Gapped (\(\ge 75\%\)) Count (\%)
\end{minipage} & \begin{minipage}[b]{\linewidth}\raggedright
Method 2: edlib Gapped (\(\ge 80\%\)) Count (\%)
\end{minipage} & \begin{minipage}[b]{\linewidth}\raggedright
Published SADDLE Figure 3f (Xie et al. 2022)
\end{minipage} & \begin{minipage}[b]{\linewidth}\raggedright
Replication Status
\end{minipage} \\
\midrule\noalign{}
\endfirsthead
\toprule\noalign{}
\begin{minipage}[b]{\linewidth}\raggedright
Read Category
\end{minipage} & \begin{minipage}[b]{\linewidth}\raggedright
Method 1: Primer Pairing Count (\%)
\end{minipage} & \begin{minipage}[b]{\linewidth}\raggedright
Method 2: edlib Gapped (\(\ge 75\%\)) Count (\%)
\end{minipage} & \begin{minipage}[b]{\linewidth}\raggedright
Method 2: edlib Gapped (\(\ge 80\%\)) Count (\%)
\end{minipage} & \begin{minipage}[b]{\linewidth}\raggedright
Published SADDLE Figure 3f (Xie et al. 2022)
\end{minipage} & \begin{minipage}[b]{\linewidth}\raggedright
Replication Status
\end{minipage} \\
\midrule\noalign{}
\endhead
\bottomrule\noalign{}
\endlastfoot
\textbf{On-Target Amplicons} & 1,254,647 (\textbf{43.86\%}) & 1,247,496
(\textbf{43.61\%}) & 1,223,216 (\textbf{42.76\%}) & \textasciitilde43\%
& \textbf{PASS} (reproduces \textasciitilde43\%) \\
\textbf{Nonspecific Background} & 1,581,139 (\textbf{55.27\%}) &
1,588,290 (\textbf{55.52\%}) & 1,612,570 (\textbf{56.37\%}) &
\textasciitilde56\% & \textbf{PASS} (reproduces \textasciitilde56\%) \\
\textbf{Primer-Dimers} & 24,872--24,878 (\textbf{0.87\%}) & 24,878
(\textbf{0.87\%}) & 24,878 (\textbf{0.87\%}) & \textasciitilde1\% &
\textbf{PASS} (invariant replicate) \\
\textbf{Total Classified Reads} & 2,860,658--2,860,664 (100\%) &
2,860,664 (100.00\%) & 2,860,664 (100.00\%) & 100\% & \textbf{PASS} \\
\end{longtable}

\subsubsection{Biological and Methodological
Insights}\label{biological-and-methodological-insights}

\begin{enumerate}
\def\labelenumi{\arabic{enumi}.}
\tightlist
\item
  \textbf{Invariance of Primer-Dimer Suppression:} Under all
  classification definitions, primer-dimer formation accounts for
  exactly \textbf{0.87\%} of classified reads, replicating the published
  finding of \textbf{\textasciitilde1\%} in Figure 3f within 0.13\%
  absolute difference. Because primer-dimers are defined strictly by
  insert length truncation (\(< L_1 + L_2 - 2\) bp), dimer
  quantification is completely robust to down-stream sequence alignment
  criteria.
\item
  \textbf{Impact of Genuine Gapped Alignment vs.~Positional String
  Matching:} In earlier exploratory audits, evaluating reference
  sequence matching via ungapped character-by-character string
  comparison artificially depressed on-target classification to 39.49\%
  (1,129,696 reads). Diagnostic inspection revealed that 1--2 bp
  sequencing indels or homopolymer stutter in MiniSeq reads caused
  downstream sequence frames to shift, artificially failing rigid
  unaligned positional checks. When genuine gapped semi-global alignment
  is applied using \texttt{edlib}, these indel-containing reads align
  with high fidelity, restoring on-target classification to
  \textbf{43.61\%} (\(\ge 75\%\) identity, \(\Delta = +0.61\%\)) and
  \textbf{42.76\%} (\(\ge 80\%\) identity, \(\Delta = -0.24\%\)). Both
  aligners quantitatively replicate the published \textasciitilde43\%
  on-target and \textasciitilde56\% nonspecific fractions within 0.6\%
  absolute tolerance.
\item
  \textbf{Turnkey Replication Concordance:} Both Method 1 and Method 2
  confirm the central empirical claim of SADDLE: multiplex PCR libraries
  designed under SADDLE's simulated annealing Badness minimization
  achieve \(\approx 43\%\) on-target specificity while rigorously
  suppressing primer-dimer amplification to under 1\%.
\end{enumerate}

\subsection{Extreme Heavy-Tailed Interaction
Concentration}\label{extreme-heavy-tailed-interaction-concentration}

We next analyzed the distribution of the observed primer-dimer reads
across all 295,296 possible candidate primer pairs: - \textbf{Zero-Dimer
Pairs:} 294,124 candidate pairs (99.6031\%) produced exactly zero
observed dimer reads. - \textbf{Active Dimer Pairs (\(\ge 1\) read):}
Only 1,172 pairs (0.3969\% empirical prevalence) generated one or more
dimer reads. - \textbf{Severe Dimer Pairs (\(\ge 10\) reads):} Only 98
pairs (0.0332\% prevalence) produced 10 or more dimer reads. -
\textbf{Dominant Interaction:} A single primer pair---reverse primer 213
(\texttt{rp213}) \(\times\) reverse primer 371
(\texttt{rp371})---generated 16,805 dimer reads, accounting for
\textbf{67.57\%} of all primer-dimer reads in the entire library.

This reproduces the qualitative finding highlighted by Xie et al.:
primer-dimer formation is not uniformly distributed across sequence
space, but is characterized by an extreme, heavy-tailed concentration
where a tiny minority of catastrophic pairs account for the vast
majority of artifactual reads.

\section{Supplementary Exploratory Analysis: Heuristic and Baseline
Benchmarking}\label{supplementary-exploratory-analysis-heuristic-and-baseline-benchmarking}

\emph{Note: The following benchmark was not published by Xie et al.
(2022). It is conducted here as a secondary exploration to evaluate how
the biophysical SADDLE Badness heuristic and machine learning models
correlate with empirical dimer labels derived from Sample E19 across all
295,296 candidate primer pairs.}

\subsection{Formulation of SADDLE Badness and
Baselines}\label{formulation-of-saddle-badness-and-baselines}

In Xie et al. (2022), Badness is calculated over maximal continuous
complementary matches of length \(\ge 4\) bp between primers \(P_1\) and
\(P_2\):

\[\text{Badness}(P_1, P_2) = \sum_{k \in \mathcal{M}} 2^{\text{len}_k} \cdot 2^{\text{numGC}_k} \cdot \frac{1}{(d_1^{(k)} + 1)(d_2^{(k)} + 1)}\]

where \(\text{len}_k\) is match length, \(\text{numGC}_k\) is GC count
within the match, and \(d_1^{(k)}, d_2^{(k)}\) are nucleotide distances
from the 3' termini. In our benchmark, we evaluate an approximate 4--7
bp heuristic variant.

We compare Badness against: 1. \textbf{SantaLucia Thermodynamics
(\(-\Delta G\)):} Hybridization free energy calculated at 60 °C using
unified nearest-neighbor parameters (SantaLucia 1998). 2.
\textbf{Logistic Regression (Physics Features):} L2-regularized linear
model trained on thermodynamic and alignment features. 3.
\textbf{XGBoost (Physics Features):} Gradient boosted decision tree
ensemble trained on the same feature set.

To ensure deterministic evaluation without tie-breaking instability,
candidates are ranked by primary score descending, with ties broken by
ascending pair index:
\(\text{rank\_key}(i) = (-\text{score}_i, \text{pair\_index}_i)\).

\subsection{Benchmark Metrics and Calibration
Analysis}\label{benchmark-metrics-and-calibration-analysis}

Table 3 presents the evaluation across all 295,296 candidate pairs
against empirical binary labels (\(\ge 1\) dimer read).

\begin{longtable}[]{@{}
  >{\raggedright\arraybackslash}p{(\columnwidth - 14\tabcolsep) * \real{0.1250}}
  >{\raggedright\arraybackslash}p{(\columnwidth - 14\tabcolsep) * \real{0.1250}}
  >{\raggedright\arraybackslash}p{(\columnwidth - 14\tabcolsep) * \real{0.1250}}
  >{\raggedright\arraybackslash}p{(\columnwidth - 14\tabcolsep) * \real{0.1250}}
  >{\raggedright\arraybackslash}p{(\columnwidth - 14\tabcolsep) * \real{0.1250}}
  >{\raggedright\arraybackslash}p{(\columnwidth - 14\tabcolsep) * \real{0.1250}}
  >{\raggedright\arraybackslash}p{(\columnwidth - 14\tabcolsep) * \real{0.1250}}
  >{\raggedright\arraybackslash}p{(\columnwidth - 14\tabcolsep) * \real{0.1250}}@{}}
\caption{Supplementary benchmark of sequence-based heuristics and
supervised models across 295,296 candidate primer pairs from the
384-plex validation panel. Binary labels represent observed dimer
formation (\(\ge 1\) NGS read). Precision@\(K\) evaluated using
deterministic secondary key sorting.}\tabularnewline
\toprule\noalign{}
\begin{minipage}[b]{\linewidth}\raggedright
Model / Heuristic
\end{minipage} & \begin{minipage}[b]{\linewidth}\raggedright
Scoring Type
\end{minipage} & \begin{minipage}[b]{\linewidth}\raggedright
ROC-AUC
\end{minipage} & \begin{minipage}[b]{\linewidth}\raggedright
PR-AUC (AP)
\end{minipage} & \begin{minipage}[b]{\linewidth}\raggedright
Precision@10
\end{minipage} & \begin{minipage}[b]{\linewidth}\raggedright
Precision@50
\end{minipage} & \begin{minipage}[b]{\linewidth}\raggedright
Precision@100
\end{minipage} & \begin{minipage}[b]{\linewidth}\raggedright
Brier Score
\end{minipage} \\
\midrule\noalign{}
\endfirsthead
\toprule\noalign{}
\begin{minipage}[b]{\linewidth}\raggedright
Model / Heuristic
\end{minipage} & \begin{minipage}[b]{\linewidth}\raggedright
Scoring Type
\end{minipage} & \begin{minipage}[b]{\linewidth}\raggedright
ROC-AUC
\end{minipage} & \begin{minipage}[b]{\linewidth}\raggedright
PR-AUC (AP)
\end{minipage} & \begin{minipage}[b]{\linewidth}\raggedright
Precision@10
\end{minipage} & \begin{minipage}[b]{\linewidth}\raggedright
Precision@50
\end{minipage} & \begin{minipage}[b]{\linewidth}\raggedright
Precision@100
\end{minipage} & \begin{minipage}[b]{\linewidth}\raggedright
Brier Score
\end{minipage} \\
\midrule\noalign{}
\endhead
\bottomrule\noalign{}
\endlastfoot
\textbf{SADDLE Badness} & Biophysical Heuristic & 0.7166 &
\textbf{0.0597} & \textbf{40.0\%} & \textbf{34.0\%} & \textbf{24.0\%} &
N/A \\
\textbf{SantaLucia (\(-\Delta G\))} & Thermodynamic (\(\Delta G\)) &
0.7105 & 0.0125 & 20.0\% & 10.0\% & 6.0\% & N/A \\
\textbf{Logistic Regression} & Supervised Probabilistic & 0.7317 &
0.0205 & 10.0\% & 4.0\% & 8.0\% & 0.1986 \\
\textbf{XGBoost} & Supervised Probabilistic & \textbf{0.7564} & 0.0371 &
20.0\% & 20.0\% & 16.0\% & 0.2143 \\
\emph{Constant Baseline} & \emph{Empirical Base Rate} & \emph{0.5000} &
\emph{0.0040} & \emph{0.4\%} & \emph{0.4\%} & \emph{0.4\%} &
\emph{0.0040} \\
\end{longtable}

\subsubsection{Findings from Supplementary
Benchmarking}\label{findings-from-supplementary-benchmarking}

\begin{enumerate}
\def\labelenumi{\arabic{enumi}.}
\tightlist
\item
  \textbf{Heuristic vs.~Thermodynamics:} SADDLE Badness achieves an
  Average Precision of 0.0597, representing a \textbf{4.8-fold
  improvement over nearest-neighbor free energy alone} (\(-\Delta G\),
  AP = 0.0125) and a \textbf{4.0-fold improvement in Precision@100}
  (24.0\% vs.~6.0\%). This supports the hypothesis of Xie et al.~that
  terminal 3' alignment and heavy GC penalties capture enzymatic
  extension risk more effectively than global hybridization stability.
\item
  \textbf{Top-Rank Discrimination vs.~Supervised Models:} SADDLE Badness
  outperforms both Logistic Regression (AP = 0.0205, P@100 = 8.0\%) and
  XGBoost (AP = 0.0371, P@100 = 16.0\%) in top-rank precision,
  demonstrating that simple biophysical heuristics remain competitive at
  the extreme head of the distribution.
\item
  \textbf{Severe Probability Miscalibration of Supervised Models:}
  Independent recalculation of Brier scores reveals that both Logistic
  Regression (0.1986) and XGBoost (0.2143) are severely miscalibrated
  relative to the uninformative baseline
  (\(p(1-p) = 0.003953 \approx 0.0040\)). Because the supervised models
  were trained to maximize ranking separation under class imbalance,
  their mean predicted probabilities (\textasciitilde0.38) greatly
  exceed the true empirical prevalence (0.00397), yielding large
  quadratic calibration penalties.
\end{enumerate}

\section{Methodological Disclosures}\label{methodological-disclosures}

\begin{enumerate}
\def\labelenumi{\arabic{enumi}.}
\tightlist
\item
  \textbf{Dual Evaluation: Primer Pairing vs.~Genuine Gapped Sequence
  Alignment:} Rather than relying solely on primer-pairing heuristics or
  unaligned string comparisons, our replication evaluated both
  classification paradigms across all 3,942,385 reads:

  \begin{itemize}
  \tightlist
  \item
    \emph{Method 1 (Primer Pairing):} Assigns on-target status based on
    cognate primer IDs (\(P_2 = P_1 + 384\)) and insert length threshold
    \(\ge L_1 + L_2 - 2\) bp, yielding 43.86\% on-target, 55.27\%
    nonspecific, and 0.87\% primer-dimers. This reproduces published
    Figure 3f (\textasciitilde43\% / \textasciitilde56\% /
    \textasciitilde1\%) within 1.0\% absolute tolerance.
  \item
    \emph{Method 2 (Genuine Gapped Sequence Alignment via edlib):}
    Performs Myers' exact bit-vector semi-global sequence alignment
    against all 384 intended reference amplicons from Supplementary Data
    MOESM5 (Xie et al. 2022). At \(\ge 75\%\) identity, reads partition
    into 43.61\% on-target, 55.52\% nonspecific, and 0.87\%
    primer-dimers (\(\Delta = +0.61\%\)). At \(\ge 80\%\) identity,
    reads partition into 42.76\% on-target, 56.37\% nonspecific, and
    0.87\% primer-dimers (\(\Delta = -0.24\%\)). Both aligner cutoffs
    match published Figure 3f within scientific tolerance.
  \item
    \emph{Positional vs Gapped Alignment Insight:} Naive ungapped string
    matching artificially depresses on-target reads to 39.49\% due to
    1--2 bp sequencing indels, whereas genuine gapped alignment robustly
    tolerates sequencing variations. Crucially, primer-dimer
    quantification (0.87\% vs \textasciitilde1\% published) is
    completely invariant across all methods.
  \end{itemize}
\item
  \textbf{Scope Limitation:} This replication focuses strictly on the
  empirical validation read partitions and heuristic scoring. The
  proprietary simulated annealing optimization algorithm remains
  withheld under patent restrictions and is not evaluated.
\item
  \textbf{Distributional Terminology:} We characterize the distribution
  of primer-dimer read counts as heavy-tailed and highly concentrated;
  no formal power-law distribution fit (e.g., Clauset-Shalizi-Newman
  methodology) was conducted.
\item
  \textbf{Data and Code Licensing:} Raw sequencing data from Figshare
  (DOI: 10.6084/m9.figshare.16944154.v3) are licensed under CC BY 4.0.
  The replication codebase is distributed under the MIT License.
\end{enumerate}

\section{Conclusion}\label{conclusion}

We have performed an independent computational replication of the
empirical NGS validation results from SADDLE (Xie et al. 2022). Using an
open-source Python workflow, we confirmed that classified reads from the
SADDLE-designed 384-plex validation panel (Figure 3f) replicate
published values (\textasciitilde43\% on-target, \textasciitilde56\%
nonspecific, \textasciitilde1\% primer-dimers). Under primer-pairing
classification (Method 1), reads partition into 43.86\% on-target,
55.27\% nonspecific, and 0.87\% primer-dimers. Under genuine gapped
sequence alignment against the 384 intended target amplicons (Method 2
via \texttt{edlib}), reads partition into 43.61\% on-target
(\(\ge 75\%\) identity) or 42.76\% on-target (\(\ge 80\%\) identity),
55.52\% to 56.37\% nonspecific, and 0.87\% primer-dimers, reproducing
the published \textasciitilde43\% / \textasciitilde56\% /
\textasciitilde1\% figures within 0.6\% absolute tolerance. Under all
criteria, primer-dimer formation is strictly suppressed at 0.87\%,
confirming the core experimental claim of SADDLE. We also corroborated
the severe heavy-tailed concentration of dimer reads, with a single
primer pair generating 67.57\% of all dimer reads. Finally, our
supplementary exploratory benchmark illustrated the predictive utility
of the SADDLE Badness heuristic over standard thermodynamics while
documenting the calibration characteristics of supervised models.

All data, scripts, and a standalone verification suite are available in
the public replication repository.

\section{Code and Data Availability}\label{code-and-data-availability}

\begin{itemize}
\tightlist
\item
  \textbf{GitHub Repository:}
  \url{https://github.com/Rashidmstar12/saddle-replication}
\item
  \textbf{Zenodo Deposit:} \url{https://doi.org/10.5281/zenodo.24100087}
\item
  \textbf{Primary Data Source:} Figshare (DOI:
  \href{https://doi.org/10.6084/m9.figshare.16944154.v3}{10.6084/m9.figshare.16944154.v3},
  CC BY 4.0)
\end{itemize}

Reviewers can verify all primary and supplementary results in under two
seconds:

\begin{Shaded}
\begin{Highlighting}[]
\FunctionTok{git}\NormalTok{ clone https://github.com/Rashidmstar12/saddle{-}replication.git}
\BuiltInTok{cd}\NormalTok{ saddle{-}replication}
\ExtensionTok{python}\NormalTok{ code/verify\_replication.py}
\end{Highlighting}
\end{Shaded}

To run end-to-end classification directly from the raw FASTQ sequencing
files:

\begin{Shaded}
\begin{Highlighting}[]
\ExtensionTok{python}\NormalTok{ code/reproduce\_from\_fastq.py }\AttributeTok{{-}{-}full}
\end{Highlighting}
\end{Shaded}

\section*{References}\label{references}
\addcontentsline{toc}{section}{References}

\phantomsection\label{refs}
\begin{CSLReferences}{1}{0}
\bibitem[\citeproctext]{ref-chen2026prism}
Chen, Yixin, Yunheng Han, Ao Wang, Aaron Hong, Adam R. Rivers, Alan
Kuhnle, and Christina Boucher. 2026. {``Primer Design Through Submodular
Function Estimation.''} \emph{Bioinformatics} 42 (7): btag478.
\url{https://doi.org/10.1093/bioinformatics/btag478}.

\bibitem[\citeproctext]{ref-langmead2012fast}
Langmead, Ben, and Steven L. Salzberg. 2012. {``Fast Gapped-Read
Alignment with {Bowtie} 2.''} \emph{Nature Methods} 9 (4): 357--59.
\url{https://doi.org/10.1038/nmeth.1923}.

\bibitem[\citeproctext]{ref-rougier2017sustainable}
Rougier, Nicolas P., Konrad Hinsen, Frédéric Oliva, Sergio Mattioni,
Karthik Ram, Patrick Vandewalle, Zheng Zou, Christopher Williams, Eric
Visser, and Li Zhang. 2017. {``Sustainable Computational Science: The
{ReScience} Initiative.''} \emph{PeerJ Computer Science} 3: e142.
\url{https://doi.org/10.7717/peerj-cs.142}.

\bibitem[\citeproctext]{ref-santalucia1998unified}
SantaLucia, John. 1998. {``A Unified View of Polymer, Dumbbell, and
Oligonucleotide {DNA} Nearest-Neighbor Thermodynamics.''}
\emph{Proceedings of the National Academy of Sciences} 95 (4): 1460--65.
\url{https://doi.org/10.1073/pnas.95.4.1460}.

\bibitem[\citeproctext]{ref-untergasser2012primer3}
Untergasser, Andreas, Ioana Cutcutache, Triinu Koressaar, Jian Ye, Brant
C. Faircloth, Maido Remm, and Steven G. Rozen. 2012. {``{Primer3}---New
Capabilities and Interfaces.''} \emph{Nucleic Acids Research} 40 (15):
e115. \url{https://doi.org/10.1093/nar/gks596}.

\bibitem[\citeproctext]{ref-xie2022designing}
Xie, Nina G., Michael X. Wang, Ping Song, Shiqi Mao, Yifan Wang, Yuxia
Yang, Junfeng Luo, Shengxiang Ren, and David Yu Zhang. 2022.
{``Designing Highly Multiplex {PCR} Primer Sets with {Simulated}
{Annealing} {Design} Using {Dimer} {Likelihood} {Estimation}
({SADDLE}).''} \emph{Nature Communications} 13 (1): 1881.
\url{https://doi.org/10.1038/s41467-022-29500-4}.

\end{CSLReferences}

\end{document}
